\subsection{Operaciones con subespacios}

\begin{frame}[allowframebreaks, t]
    \frametitle{Intersección de subespacios}

    \begin{proposition}
        Si $W_1$ y $W_2$ son subespacios de $V$, entonces
        $W_1\cap W_2$ también lo es.
    \end{proposition}

    \begin{proof}
        Como $0\in W_1$ y $0\in W_2$, se tiene $0\in W_1\cap W_2$.
        Sean $u,v\in W_1\cap W_2$ y $\lambda\in k$.
        Como $W_1$ es un subespacio, $u+v\in W_1$ y $\lambda u\in W_1$;
        análogamente para $W_2$.
        Por tanto $u+v\in W_1\cap W_2$ y $\lambda u\in W_1\cap W_2$.
    \end{proof}
\end{frame}


\begin{frame}
    \frametitle{Ecuaciones de la intersección en $k^n$}

    \begin{proposition}
        Sean $A_1\in M_{m_1\times n}(k)$ y $A_2\in M_{m_2\times n}(k)$. Si
        $W_1=\{v\in k^n:\ A_1v=0\}$ y $W_2=\{v\in k^n:\ A_2v=0\}$,
        entonces
        \[
            W_1\cap W_2=\left\{v\in k^n:\
            \begin{pmatrix}A_1\\A_2\end{pmatrix}v=0\right\}.
        \]
        Es decir, la unión de ecuaciones implícitas de $W_1$ y $W_2$
        es un sistema de ecuaciones implícitas de $W_1\cap W_2$.
    \end{proposition}

    \begin{proof}
        Las primeras $m_1$ entradas de
        $\begin{pmatrix}A_1\\A_2\end{pmatrix}v$ son las de $A_1v$, y las
        $m_2$ restantes son las de $A_2v$. Por tanto, este vector es nulo
        si y solo si $A_1v=0$ y $A_2v=0$, es decir, si y solo si
        $v\in W_1$ y $v\in W_2$.
    \end{proof}
\end{frame}


\begin{frame}
    \frametitle{Ejemplo: intersección de dos planos en $\R^3$}

    \begin{example}
        Sean $W_1,W_2\subseteq\R^3$ los planos
        \[
            W_1\colon\ v_1-2v_2+v_3=0,
            \qquad
            W_2\colon\ v_1+v_2=0.
        \]
        Calculemos $W_1\cap W_2$.
    \end{example}

    Por la proposición anterior, $W_1\cap W_2$ es el conjunto de
    soluciones del sistema formado por las dos ecuaciones. Estas son
    independientes (no son proporcionales), luego
    $\dim(W_1\cap W_2)=3-2=1$. De la segunda, $v_1=-v_2$; sustituyendo en
    la primera, $v_3=3v_2$. Por tanto,
    \[
        W_1\cap W_2=\{(-v_2,v_2,3v_2):\ v_2\in\R\}=\langle(-1,1,3)\rangle,
    \]
    una recta que pasa por el origen. En general, la unión de ecuaciones
    independientes de $W_1$ y de $W_2$ \alert{no} tiene por qué ser
    independiente; si no lo es, se escalona para eliminar las redundantes.
\end{frame}


\begin{frame}
    \frametitle{Ecuaciones de la intersección respecto de una base}

    \begin{proposition}
        Sea $V$ un espacio vectorial de dimensión $n$ y sea $B$ una base
        de $V$. Si $W_1\colon A_1x=0$ resp.\ $B$ y
        $W_2\colon A_2x=0$ resp.\ $B$, entonces
        \[
            W_1\cap W_2\colon\
            \begin{pmatrix}A_1\\A_2\end{pmatrix}x=0\quad\text{resp.\ }B.
        \]
    \end{proposition}

    \begin{proof}
        Para $u\in V$, aplicamos la proposición anterior a $v=u_B$:
        \[
            u\in W_1\cap W_2
            \iff A_1u_B=0\ \text{y}\ A_2u_B=0
            \iff \begin{pmatrix}A_1\\A_2\end{pmatrix}u_B=0.\qedhere
        \]
    \end{proof}

    Las ecuaciones de $W_1$ y $W_2$ deben estar dadas respecto de la
    \alert{misma base}; si no, antes hay que cambiar de base una de ellas.
\end{frame}


\begin{frame}
    \frametitle{Ejemplo: intersección de subespacios de matrices}

    \begin{example}
        En $M_2(k)$, sean $W_1$ el subespacio de las matrices simétricas
        y $W_2$ el de las triangulares superiores. Calculemos
        $W_1\cap W_2$.
    \end{example}

    Respecto de la base
    \[
        B=\left\{
        \begin{pmatrix}1&0\\0&0\end{pmatrix},
        \begin{pmatrix}0&1\\0&0\end{pmatrix},
        \begin{pmatrix}0&0\\1&0\end{pmatrix},
        \begin{pmatrix}0&0\\0&1\end{pmatrix}
        \right\}
    \]
    ya vimos que $W_1\colon x_2-x_3=0$ y $W_2\colon x_3=0$. Por tanto,
    \[
        W_1\cap W_2\colon\
        \begin{pmatrix}0&1&-1&0\\0&0&1&0\end{pmatrix}x=0
        \quad\text{resp.\ }B,
    \]
    es decir, $x_2=x_3=0$: la intersección está formada por las matrices
    \alert{diagonales}
    $\begin{pmatrix}a&0\\0&d\end{pmatrix}$.
\end{frame}


\begin{frame}
    \frametitle{La unión de subespacios no es un subespacio}

    \begin{columns}[T]
        \begin{column}{0.62\textwidth}
            A diferencia de la intersección, la unión $W_1\cup W_2$
            no es en general un subespacio.

            \medskip

            Ya vimos que la unión de los dos ejes de $\R^2$ no lo es:
            $(1,0)$ y $(0,1)$ están en ella, pero
            \[
                (1,0)+(0,1)=(1,1)
            \]
            no está en ninguno de los dos ejes.
        \end{column}
        \begin{column}{0.34\textwidth}
            \centering
            \begin{tikzpicture}[scale=0.85]
                \draw[very thick,alertcolor] (-1.6,0) -- (2,0);
                \draw[very thick,alertcolor] (0,-1.2) -- (0,2);
                \node[alertcolor,right] at (2,0) {$W_1$};
                \node[alertcolor,above] at (0,2) {$W_2$};
                \draw[dashed] (1,0) -- (1,1) -- (0,1);
                \fill[alertcolor] (1,0) circle (2.5pt)
                    node[below] {\footnotesize $(1,0)$};
                \fill[alertcolor] (0,1) circle (2.5pt)
                    node[left] {\footnotesize $(0,1)$};
                \fill[white] (1,1) circle (2.5pt);
                \draw[thick] (1,1) circle (2.5pt)
                    node[above right] {\footnotesize $(1,1)$};
            \end{tikzpicture}
        \end{column}
    \end{columns}

    \medskip

    El menor subespacio que contiene a $W_1$ y a $W_2$ es la
    \alert{suma}, que definimos a continuación.
\end{frame}


\begin{frame}[allowframebreaks, t]
    \frametitle{Suma de subespacios}

    \begin{definition}
        La \alert{suma} de dos subespacios $W_1$ y $W_2$ de $V$ es
        \[
            W_1+W_2=\{u+w:\; u\in W_1,\; w\in W_2\}.
        \]
    \end{definition}

    \begin{proposition}
        $W_1+W_2$ es el menor subespacio de $V$ que contiene a $W_1$
        y a $W_2$.
    \end{proposition}

    \begin{proof}
        Es un subespacio: $0=0+0\in W_1+W_2$; si $u_1+w_1,u_2+w_2\in
        W_1+W_2$ y $\lambda\in k$, entonces
        \[
            (u_1+w_1)+(u_2+w_2)=(u_1+u_2)+(w_1+w_2)\in W_1+W_2,
        \]
        \[
            \lambda(u_1+w_1)=\lambda u_1+\lambda w_1\in W_1+W_2.
        \]
        Contiene a $W_1$ (tomando $w=0$) y a $W_2$ (tomando $u=0$).
        Si $W$ es cualquier subespacio que contiene a $W_1$ y $W_2$,
        entonces $u+w\in W$ para todo $u\in W_1$, $w\in W_2$,
        luego $W_1+W_2\subseteq W$.
    \end{proof}
\end{frame}


\begin{frame}[allowframebreaks, t]
    \frametitle{Generadores de la suma}

    \begin{proposition}
        Si $W_1=\langle v_1,\ldots,v_p\rangle$ y
        $W_2=\langle w_1,\ldots,w_q\rangle$, entonces
        \[
            W_1+W_2=\langle v_1,\ldots,v_p,w_1,\ldots,w_q\rangle.
        \]
        Es decir, la unión de sistemas de generadores de $W_1$ y $W_2$
        es un sistema de generadores de $W_1+W_2$.
    \end{proposition}

    \begin{proof}
        Probamos las dos inclusiones.

        ($\supseteq$) Los vectores $v_1,\ldots,v_p$ están en $W_1$ y
        $w_1,\ldots,w_q$ están en $W_2$, luego todos están en $W_1+W_2$.
        Como $W_1+W_2$ es un subespacio y
        $\langle v_1,\ldots,v_p,w_1,\ldots,w_q\rangle$ es el menor
        subespacio que contiene a estos vectores, se tiene
        $\langle v_1,\ldots,v_p,w_1,\ldots,w_q\rangle\subseteq W_1+W_2$.

        ($\subseteq$) Todo elemento de $W_1+W_2$ es de la forma $u+w$ con
        $u\in W_1$ y $w\in W_2$. Como $v_1,\ldots,v_p$ generan $W_1$ y
        $w_1,\ldots,w_q$ generan $W_2$, existen escalares
        $\lambda_1,\ldots,\lambda_p,\mu_1,\ldots,\mu_q\in k$ tales que
        \[
            u=\lambda_1v_1+\cdots+\lambda_pv_p,
            \qquad
            w=\mu_1w_1+\cdots+\mu_qw_q.
        \]
        Sumando,
        \[
            u+w=\lambda_1v_1+\cdots+\lambda_pv_p+\mu_1w_1+\cdots+\mu_qw_q,
        \]
        que es una combinación lineal de
        $v_1,\ldots,v_p,w_1,\ldots,w_q$.
        Por tanto $W_1+W_2\subseteq\langle v_1,\ldots,v_p,w_1,\ldots,w_q\rangle$.
    \end{proof}
\end{frame}


\begin{frame}
    \frametitle{La unión de bases no es una base de la suma}

    \begin{example}
        En $\R^3$, sean
        \[
            W_1=\langle(1,0,0),(0,1,0)\rangle,
            \qquad
            W_2=\langle(1,1,0),(0,0,1)\rangle.
        \]
        Los dos pares de generadores son linealmente independientes, luego
        son bases de $W_1$ y $W_2$. Su unión genera $W_1+W_2$, pero
        \alert{no} es una base de $W_1+W_2$.
    \end{example}

    Por la proposición anterior,
    \[
        W_1+W_2=\langle(1,0,0),(0,1,0),(1,1,0),(0,0,1)\rangle=\R^3,
    \]
    pues entre estos vectores están $(1,0,0)$, $(0,1,0)$ y $(0,0,1)$,
    que generan $\R^3$.
\end{frame}


\begin{frame}
    \frametitle{La unión de bases no es una base de la suma}

    Sin embargo, los cuatro vectores
    \[
        (1,0,0),\quad(0,1,0),\quad(1,1,0),\quad(0,0,1)
    \]
    son linealmente dependientes:
    \[
        1\cdot(1,0,0)+1\cdot(0,1,0)+(-1)\cdot(1,1,0)+0\cdot(0,0,1)=0,
    \]
    con coeficientes no todos nulos.

    \medskip

    El problema es que $(1,1,0)$ está en $W_1\cap W_2$: es a la vez
    combinación lineal de la base de $W_1$ y elemento de la base de $W_2$.
    La unión de bases genera la suma, pero en general
    \alert{no es independiente}. Para obtener una base de $W_1+W_2$ hay que
    eliminar vectores sobrantes; aquí basta quitar $(1,1,0)$.
\end{frame}


\begin{frame}
    \frametitle{Ejemplo: intersección y suma en $\R^3$}

    \begin{example}
        En $\R^3$, sean $W_1=\{u_3=0\}=\langle e_1,e_2\rangle$ y
        $W_2=\{u_1=0\}=\langle e_2,e_3\rangle$. Entonces
        \[
            W_1\cap W_2=\{u_1=0,\;u_3=0\}=\langle e_2\rangle,
            \qquad
            W_1+W_2=\R^3.
        \]
    \end{example}

    Para $W_1+W_2$: todo vector $(u_1,u_2,u_3)\in\R^3$ se descompone como
    \[
        (u_1,u_2,u_3)=\underbrace{(u_1,u_2,0)}_{\in W_1}
        +\underbrace{(0,0,u_3)}_{\in W_2}.
    \]
\end{frame}


\begin{frame}[allowframebreaks, t]
    \frametitle{Fórmula de Grassmann}

    \begin{theorem}
        Si $W_1$ y $W_2$ son subespacios de dimensión finita de $V$, entonces
        \[
            \dim(W_1+W_2)=\dim W_1+\dim W_2-\dim(W_1\cap W_2).
        \]
    \end{theorem}

    \begin{proof}
        Sea $\{u_1,\ldots,u_d\}$ una base de $W_1\cap W_2$, que
        extendemos a bases
        \[
            \{u_1,\ldots,u_d,v_1,\ldots,v_p\}\quad\text{de }W_1,
            \qquad
            \{u_1,\ldots,u_d,w_1,\ldots,w_q\}\quad\text{de }W_2.
        \]
        Afirmamos que $B=\{u_1,\ldots,u_d,v_1,\ldots,v_p,w_1,\ldots,w_q\}$
        es una base de $W_1+W_2$.

        \alert{Generación.}
        $B$ es la unión de una base de $W_1$ y una base de $W_2$,
        luego genera $W_1+W_2$ por la proposición sobre generadores
        de la suma.

        \alert{Independencia lineal.}
        Supongamos $\sum a_i u_i+\sum b_j v_j+\sum c_k w_k=0$.
        Entonces
        \[
            \sum c_k w_k=-\sum a_i u_i-\sum b_j v_j\in W_1.
        \]
        Como $\sum c_k w_k\in W_2$, se tiene $\sum c_k w_k\in W_1\cap W_2$,
        luego $\sum c_k w_k=\sum\lambda_i u_i$ para ciertos $\lambda_i$.
        Reordenando, $-\sum\lambda_i u_i+\sum c_k w_k=0$.
        Como $\{u_1,\ldots,u_d,w_1,\ldots,w_q\}$ es base de $W_2$,
        concluimos $c_k=0$ para todo $k$.
        Entonces $\sum a_i u_i+\sum b_j v_j=0$, y como
        $\{u_1,\ldots,u_d,v_1,\ldots,v_p\}$ es base de $W_1$,
        también $a_i=b_j=0$.

        Por tanto $\dim(W_1+W_2)=d+p+q=\dim W_1+\dim W_2-d$.
    \end{proof}
\end{frame}


\begin{frame}
    \frametitle{Consecuencias de la fórmula de Grassmann}

    Para $W_1=\langle e_1,e_2\rangle$ y $W_2=\langle e_2,e_3\rangle$ en $\R^3$:
    \[
        \dim(W_1+W_2)=2+2-1=3,
    \]
    lo que confirma $W_1+W_2=\R^3$.

    \begin{corollary}
        Si $\dim V=n$ y $W_1,W_2\subseteq V$, entonces
        \[
            \dim(W_1\cap W_2)\geq\dim W_1+\dim W_2-n.
        \]
        En particular, si $\dim W_1+\dim W_2>n$, la intersección
        no puede ser trivial.
    \end{corollary}
\end{frame}


\begin{frame}[allowframebreaks, t]
    \frametitle{Suma directa}

    \begin{definition}
        Se dice que $V$ es la \alert{suma directa} de $W_1$ y $W_2$,
        y se escribe $V=W_1\oplus W_2$, si
        $V=W_1+W_2$ y $W_1\cap W_2=\{0\}$.
    \end{definition}

    \begin{proposition}
        $V=W_1\oplus W_2$ si y solo si todo $v\in V$ se escribe de
        manera \alert{única} como $v=u+w$ con $u\in W_1$ y $w\in W_2$.
    \end{proposition}

    \begin{proof}
        ($\Rightarrow$) $V=W_1+W_2$ garantiza la existencia. Si
        $u_1+w_1=u_2+w_2$ con $u_i\in W_1$ y $w_i\in W_2$, entonces
        $u_1-u_2=w_2-w_1\in W_1\cap W_2=\{0\}$, luego $u_1=u_2$ y $w_1=w_2$.

        ($\Leftarrow$) La existencia da $V=W_1+W_2$. Si $v\in W_1\cap W_2$,
        entonces $v=v+0=0+v$ son dos descomposiciones de $v$; por
        unicidad, $v=0$.
    \end{proof}

    \begin{corollary}
        Si $V=W_1\oplus W_2$, entonces $\dim V=\dim W_1+\dim W_2$.
    \end{corollary}
\end{frame}


\begin{frame}
    \frametitle{Ejemplos de suma directa}

    \begin{example}
        En $\R^3$:
        \begin{enumerate}
            \item $\R^3=\langle e_1,e_2\rangle\oplus\langle e_3\rangle$:
            todo $(u_1,u_2,u_3)=(u_1,u_2,0)+(0,0,u_3)$, de forma única.
            \item $\R^3=\langle e_1\rangle\oplus\langle e_2\rangle\oplus\langle e_3\rangle$.
            \item $\langle e_1,e_2\rangle$ y $\langle e_2,e_3\rangle$
            no dan suma directa, pues su intersección es
            $\langle e_2\rangle\neq\{0\}$.
        \end{enumerate}
    \end{example}

    \medskip

    En general, si $\{v_1,\ldots,v_n\}$ es una base de $V$, entonces
    \[
        V=\langle v_1\rangle\oplus\cdots\oplus\langle v_n\rangle.
    \]
\end{frame}
